% Tradução da versão provisória do Resumo — atualizar junto com ele.
Large language models are usually integrated into web applications through
direct calls to the provider's software development kit, scattered across the
business code, which couples the application to a specific vendor. An
alternative is to place an intermediate architectural layer --- here called a
\textit{harness} --- that centralizes communication with providers, tool
execution, state, failure handling, and event logging. The literature proposes
such architectures and discusses their trade-offs qualitatively, but does not
measure them in controlled studies. This work designs and experimentally
evaluates this layer, weighing provider replacement cost against its energy
cost by comparing two functionally equivalent versions of the same web
application: one integrated with the provider through the harness and one integrated
directly. Provider replacement cost, understood as the code change effort
required to switch providers while preserving the application's behavior, is
quantified through behavioral and structural code metrics when switching between
the Anthropic and OpenAI providers, and energy cost through GPU hardware
telemetry for simple and agentic requests to the open model Qwen3.6 27B, hosted
on an institutional server. The work is expected to
characterize under which conditions the decoupling achieved outweighs the
energy cost of the layer, providing quantitative evidence for an architectural
decision currently made on the basis of design principles or vendor claims.

% Separe as Keywords por ponto
\keywords{Software architecture. LLM integration. Substitutability. Energy efficiency. Green AI.}
